\documentclass[11pt]{article}
\usepackage[utf8]{inputenc}
\usepackage[T1]{fontenc}
\usepackage{lmodern}
\usepackage[margin=1in]{geometry}
\usepackage{graphicx}
\usepackage{booktabs}
\usepackage{amsmath}
\usepackage{siunitx}
\usepackage{hyperref}
\usepackage{xcolor}
\usepackage{authblk}
\hypersetup{colorlinks=true, linkcolor=blue!50!black, citecolor=blue!50!black, urlcolor=blue!50!black}
\usepackage[numbers,sort&compress]{natbib}

\title{\textbf{When does wearable multi-signal fusion stop working?\\
A cross-session benchmark for fast biometric authentication}}
\author{Noor Abdalla}
\affil{\small Independent researcher \quad\textbullet\quad
\texttt{\href{https://github.com/noorabdalla04/wearable-fusion-auth-benchmark}{github.com/noorabdalla04/wearable-fusion-auth-benchmark}}}
\date{2026}

\begin{document}
\maketitle

\begin{abstract}
\noindent
Consumer wearables can read several physiological signals at once: photoplethysmography
(PPG), electrocardiography (ECG), electrodermal activity (GSR), and accelerometry (ACC). A
recurring proposal is to fuse them into a fast ``your body is your password'' unlock:
glance at the watch, get identified in seconds, like a fingerprint. Published fusion systems
report near-perfect accuracy (equal error rates around \SI{2}{\percent}), which makes the
idea look solved. We show that this number is an artefact of evaluating enrolment and test
data from the \emph{same recording session}. We build an open, data-source-agnostic benchmark
that (i) under a leak-free evaluation reaches a comparably low within-session number on a
public multi-signal dataset (Blasco et al.\ 2018; full-fusion EER \num{0.030}, in the range of
the reported ${\sim}0.02$), and
then (ii) re-evaluates the identical pipeline across sessions using two protocols: a
cross-activity proxy on the same day, and a genuine cross-day protocol on a second dataset
(PhysioNet Exam-Stress; three exams on different days). The within-session strength does not
transfer: on Blasco the best honest cross-activity EER is \num{0.199} (about a $7\times$
degradation), and on the true cross-day wrist dataset the signal is already weak within-session
(EER \num{0.234}) and stays weak across days (best \num{0.312}). Signal fusion gives a small,
real edge over the best single signal ($+0.04$ to $+0.06$ EER) but does not close the gap, and
on both datasets a cardiac channel is dropped from the cross-session optimum because its
within-session identity content does not transfer. Translated to a security operating point, the best
honest cross-session fusion is ${\sim}$\num{10000} to \num{16000}$\times$ more permissive than
Touch ID and ${\sim}$\num{200000} to \num{312000}$\times$ more permissive than Face ID. We
conclude that fast wearable fusion is not a device-unlock primitive; its honest role is a
low-stakes ``probably still the same wearer'' continuity check. All code, loaders, and a
one-command reproduction are released.
\end{abstract}

\section{Introduction}
\label{sec:intro}

Every modern smartwatch and fitness band carries an optical pulse sensor, an accelerometer,
and increasingly an electrodermal or single-lead ECG sensor. Because these signals reflect
individual physiology, a long line of work has asked whether they can \emph{identify} the
wearer, and a particularly attractive framing is authentication: use the body's own signals
as a password, so a device unlocks the moment it is worn, with no PIN and no deliberate
gesture. The appeal is obvious, since it would be effortless and continuous, and the
literature appears to support it, with multi-signal fusion systems reporting equal error
rates (EERs) around \SI{2}{\percent} or better \citep{blasco2018}.

There is a catch that is easy to miss and decisive in practice. Most flattering numbers are
obtained by splitting a \emph{single recording session} into enrolment and test partitions.
Within one session, a machine-learned classifier can lock onto session-specific nuisance
structure (sensor placement, contact pressure, skin state, the ambient conditions of that
particular sitting) that is perfectly predictive of ``which person'' \emph{within that
session} but has nothing to do with stable identity. A real unlock must work tomorrow, on a
freshly donned device, in a different place. The scientific question is therefore not ``can
we separate people within a session'' (we can, easily) but ``does that separation survive a
change of session,'' which is the only regime a password lives in.

This paper answers that question empirically and honestly. Our contributions are:
\begin{enumerate}
  \item An open, \textbf{data-source-agnostic benchmark} for fast wearable biometric
  authentication: a common recording schema plus thin per-dataset loaders, so heterogeneous
  datasets (different devices, signals, and sampling rates) are evaluated by one leakage-audited
  pipeline (\S\ref{sec:methods}).
  \item A \textbf{comparably low within-session number under a leak-free split}
  (full-fusion EER \num{0.030} on Blasco 2018, in the range of the published ${\sim}0.02$),
  establishing that our pipeline is competitive with prior art rather than a weak straw man
  (\S\ref{sec:within}).
  \item A \textbf{cross-session evaluation} under two protocols (a same-day cross-activity
  proxy and a true cross-day protocol). On Blasco, within-session fusion degrades about
  $7\times$ across activities; on the true cross-day wrist dataset the signal is weak
  within-session and stays weak across days (\S\ref{sec:cross}).
  \item A \textbf{fusion verdict and security-bar accounting}: fusion helps a little but does
  not close the gap, and the honest operating point is $4$-$5$ orders of magnitude short of
  Touch ID / Face ID (\S\ref{sec:fusion}, \S\ref{sec:security}).
\end{enumerate}

This is, deliberately, a negative result with a constructive framing. We are not reporting
that a new system works; we are measuring precisely how and where a widely-assumed idea fails,
and releasing the tooling so the measurement is cheap to repeat and hard to fudge.

\section{Related work}
\label{sec:related}

\paragraph{Physiological biometrics.} PPG and ECG have both been studied as biometric traits,
with within-session identification and verification accuracies frequently reported in the
high-90s percent \citep{blasco2018}. Fusion of cardiac signals with other modalities
(electrodermal activity, accelerometry, or a second cardiac channel) is a recurring theme,
motivated by the intuition that jointly the signals are more unique than any one alone.

\paragraph{The permanence problem.} That intuition is rarely tested across time. The
permanence of a biometric, its stability across sessions and days, is a named property
in the biometrics literature, and several wrist-PPG studies that do test it report large
degradation across days. Crucially, the dataset we anchor on states its own limitation
plainly: its signals ``were acquired only once on a given day,'' and it explicitly recommends
``cross-day and long-term analysis'' as future work \citep{blasco2018}. Our study is, in part,
that recommended follow-up, extended to fusion and to a security-bar comparison.

\paragraph{Where this work sits.} We do not claim that fusing wearable signals to identify a
person is a new idea; it is not. Nor do we claim the within-session-inflation critique is new:
a contemporaneous benchmark, ECG-biometrics-bench \citep{ecgbench2026}, names the same
``random-split'' fallacy and evaluates it across seven ECG datasets with both closed- and
open-set protocols, and is more thorough than this work on the single-signal cardiac case. Our
contribution is the complementary \emph{multi-signal fusion} slice that that benchmark does not
cover: a small, honest, reproducible study asking specifically whether \emph{fusing}
simultaneous wearable signals, the mechanism most often proposed to reach a fast unlock,
survives across sessions. We (a) reach a competitive within-session number under a leak-free
split with the same pipeline we then stress-test, (b) report cross-session numbers with
explicit leakage controls and confidence intervals, and (c) state the security implication in
practitioner units. The honest verdict is that fusion does not close the gap.

\section{Methods}
\label{sec:methods}

\subsection{Datasets}
We use two public datasets with complementary signal coverage (Table~\ref{tab:datasets}).

\paragraph{Blasco 2018 \citep{blasco2018}.} A low-cost multi-signal wearable dataset:
25 subjects (ages 18-42, mean 28.1; 16 male, 9 female), each recorded in three activity
states on a single day (seated rest; walking; seated after a short stroll). Signals are PPG,
ECG, and GSR at \SI{100}{\hertz}. There is \emph{no usable accelerometer} in the released
files (the nominal motion channels are constant or zero), so Blasco fusion here is PPG+ECG+GSR.
Because all recordings are from one day, its cross-``session'' axis is really
\emph{cross-activity}, a proxy for, not a substitute for, cross-day.

\paragraph{PhysioNet Exam-Stress.} An Empatica E4 wrist dataset: 10 subjects, each recorded
during three exams held on \emph{different days} (two midterms and a final). Signals are BVP
(PPG) at \SI{64}{\hertz}, 3-axis ACC at \SI{32}{\hertz}, and EDA (GSR) at \SI{4}{\hertz};
there is \emph{no ECG} (a wrist device). Exam-Stress fusion is therefore PPG+ACC+GSR, and its
cross-session axis is a genuine \emph{cross-day} test. Recordings span 3 to 7 hours (they
include pre/post-exam wear).

Together the two datasets cover all four signals (PPG, ECG, GSR, ACC), and provide both a
same-day proxy and a true cross-day protocol.

\begin{table}[t]
\centering
\small
\begin{tabular}{lccccc}
\toprule
Dataset & Subjects & Sessions & Signals & Rate & Cross-session axis \\
\midrule
Blasco 2018 & 25 & 3 activities (1 day) & PPG, ECG, GSR & \SI{100}{\hertz} & cross-activity (proxy) \\
Exam-Stress & 10 & 3 exams (diff.\ days) & PPG, ACC, GSR & \SI{64}{\hertz}$^\dagger$ & cross-day (true) \\
\bottomrule
\end{tabular}
\caption{The two benchmark datasets. $^\dagger$Exam-Stress channels have different native
rates (BVP 64, ACC 32, EDA \SI{4}{\hertz}); the loader resamples all channels onto a shared
\SI{64}{\hertz} grid so the schema's equal-length, time-aligned requirement holds.}
\label{tab:datasets}
\end{table}

\subsection{Common schema and data-source-agnostic loading}
A single \texttt{Recording} type holds \{subject id, session id, condition, sampling rate,
and a dictionary of equal-length time-aligned channels\}, with canonical channel names
(PPG, ECG, GSR, ACC, \dots). Each dataset has a thin loader that maps its raw files into this
type; no downstream code knows which dataset it is processing. This is what let us add a
second dataset (different device, different signals, different rates) with zero change to the
signal-processing, feature, evaluation, or analysis code. The \texttt{session\_id} field is the
split axis: within-session vs.\ cross-session evaluation differ only in how sessions are
paired, never in the pipeline.

\subsection{Signal processing and windowing}
Signals are cleaned with standard filters (NeuroKit2 for PPG/ECG; a \SI{5}{\hertz} low-pass for
GSR/EDA; ACC used as vector magnitude). We segment each recording into non-overlapping
\SI{5}{\second} windows, the fast, ``glance-and-go'' regime that a password-like unlock would
occupy. Per-window quality flags mark non-finite, flat, or saturated channels; a window is kept
if at least one cardiac channel (PPG or ECG) is good, and per-channel ``ok'' flags let the
evaluation filter per signal combination (this is what prevents a dead GSR electrode on a few
subjects from silently dropping them entirely).

\subsection{Features}
From each window we extract \emph{signal-shape} features only, never any identity- or
recording-level token. For cardiac channels: heart-rate and rate-variability statistics from
detected beats (mean HR, IBI standard deviation, RMSSD), pulse-morphology statistics
(amplitude, width, upslope), spectral band energy, and distributional shape (skew, kurtosis).
For GSR: level, spread, range, and slope. For ACC: level, spread, range, and low/high band
energy. This yields a 26-feature matrix per dataset.

\subsection{Leakage audit}
Before any evaluation we run an automated audit that (i) rejects identity-like feature names
by token-boundary matching, (ii) confirms the window index is never a feature, (iii) flags
zero-variance or degenerate features, and (iv) red-flags any single feature whose
between-subject separability is implausibly high relative to the rest (a recording-id artefact
signature). The audit passes on both datasets. As an end-to-end control we also permute
subject labels within each session: a leak-free pipeline must then score at chance, which it
does (\S\ref{sec:robust}).

\subsection{Evaluation protocol}
We evaluate closed-set verification and report the equal error rate (EER, where false-accept
rate equals false-reject rate) and the ROC AUC. The scorer is a per-window RandomForest subject
classifier (200 trees); the verification score for a claimed identity is the classifier's
predicted probability of that subject. We deliberately use a supervised per-window classifier
because that is the class of model that produces the flattering literature numbers: a
weaker distance-template scorer does \emph{not} reproduce them (it gives within-session
EER ${\approx}0.17$), so anchoring on the strong model is the honest, hardest test of our
thesis. Standardisation statistics are fit on enrolment windows only.

\begin{itemize}
  \item \textbf{Within-session:} enrol and test on disjoint window partitions of the
  \emph{same} session; average over sessions. Window partitions are asserted disjoint at
  runtime.
  \item \textbf{Cross-session:} enrol on one session, test on a \emph{different} session
  (only subjects present in both); average over ordered session pairs. On Blasco this is
  cross-activity; on Exam-Stress it is cross-day.
\end{itemize}

All randomness is seeded (\texttt{seed} in \texttt{config.yaml}); the entire pipeline is
regenerated by a single \texttt{reproduce.py}.

\section{Results}
\label{sec:results}

\subsection{A strong within-session number, under a leak-free split}
\label{sec:within}
On Blasco, full-fusion (PPG+ECG+GSR) within-session verification reaches
\textbf{EER \num{0.030}} under a leak-free time-blocked split (Table~\ref{tab:headline}), in
the range of the ${\sim}0.02$ reported by \citet{blasco2018}. We stress that this is a
comparably low number, not a like-for-like reproduction: \citet{blasco2018} used a different
scorer and 60-second training, whereas we use a per-window RandomForest. It is enough to
establish that our pipeline is competitive with the optimistic result that motivates the
``body-as-password'' premise, so the cross-session failure below is not an artefact of a weak
model. A prior version of this pipeline used a random-window split and reported EER
\num{0.014}; that split leaks temporally adjacent windows across the enrol/probe boundary, and
the leak-free number (\num{0.030}) is the honest one (see \S\ref{sec:robust}). Every single
signal is also individually usable within session (EER \num{0.113}-\num{0.212}).

\subsection{Accuracy collapses across sessions}
\label{sec:cross}
The same pipeline, evaluated across sessions, fails in two distinct ways
(Figure~\ref{fig:collapse}, Table~\ref{tab:headline}). On Blasco the best cross-activity EER
is \num{0.199}, a same-combination degradation of about $7\times$ from the within-session
\num{0.030}: a strong same-session model that does not transfer even across activity states on
the same day. On Exam-Stress, the true cross-day test on a real wrist device, the story is
different and more consequential: the signal is already weak within-session (full-fusion EER
\num{0.234}, best combination \num{0.226}) and only degrades a further ${\sim}1.5\times$ to
\num{0.312} across days. There is no strong number to collapse from, because the wrist signal
never authenticated well even within a session. At these cross-session operating points the
system is close to a coin flip at any usable threshold.

\begin{table}[t]
\centering
\small
\begin{tabular}{llccc}
\toprule
Dataset & Signals & Within EER & Cross EER & Same-combo collapse \\
\midrule
Blasco 2018 & PPG+ECG+GSR & \textbf{0.030} & 0.199-0.210 & $\sim$6.9$\times$ \\
Exam-Stress & PPG+GSR+ACC & 0.234 & 0.312-0.339 & $\sim$1.5$\times$ \\
\bottomrule
\end{tabular}
\caption{Headline result (leak-free time-blocked within-session split). On Blasco a strong
within-session fusion does not survive a change of activity state; on the true cross-day wrist
dataset the signal is weak within-session and stays weak across days. Within EER is the
full-fusion triple; cross EER ranges span the best cross-session combination and the
full-triple fusion; collapse is same-combination (full triple). Blasco cross is cross-activity
(same day); Exam-Stress cross is cross-day (true).}
\label{tab:headline}
\end{table}

\begin{figure}[t]
\centering
\includegraphics[width=\textwidth]{{{artifact:art_fd5c405b-9334-494f-8a32-9013796f9849}}}
\caption{Within-session (blue) vs.\ cross-session (red) EER for every signal combination, on
both datasets. Within-session points cluster near zero on Blasco; every combination jumps
right under a session change. Lower is better.}
\label{fig:collapse}
\end{figure}

\subsection{Does fusion survive? A small edge, not a rescue}
\label{sec:fusion}
Fusion's cross-session advantage over the best single signal is real but small: $+0.062$ EER
on Blasco, $+0.041$ on Exam-Stress. It does not close the gap; absolute cross-session EER
stays at \num{0.20}-\num{0.34} regardless. A telling detail: the full-fusion
triple is never the cross-session optimum. The best cross-session combination is PPG+GSR on
Blasco (ECG drops out) and GSR+ACC on Exam-Stress (PPG drops out): in each dataset a
cardiac channel is discarded because its within-session identity content does not transfer.
Per-signal analysis (Figure~\ref{fig:persignal}) shows that no single signal survives the move
across sessions. For a would-be unlock this is a concerning failure mode: on the wrist device
(Exam-Stress), the cardiac PPG, the signal a consumer unlock would most naturally rely on,
is the least useful single signal across days in absolute terms (cross-day EER \num{0.467},
near chance) and is dropped from the optimum entirely. Two honesty notes on this ``fusion
helps'' claim: the best cross-session combination is selected post-hoc as the minimum over
seven combinations on the same evaluation data (${\sim}6$ session-pair units), so it is
optimistically biased; and on Exam-Stress the full three-signal fusion (\num{0.339}) is
actually worse than the two-signal GSR+ACC (\num{0.312}), i.e.\ adding PPG hurts. The edge is
real but small and combination-specific, not a ``fuse everything and it improves'' result.

\begin{figure}[t]
\centering
\includegraphics[width=\textwidth]{{{artifact:art_f14625cb-e4af-4872-af62-93f43ed3aa88}}}
\caption{Per-signal within- vs.\ cross-session EER. Every signal degrades toward chance
(0.50) under a session change. On the wrist device (Exam-Stress, right) the cardiac PPG is the
least stable single signal across days and is dropped from that dataset's cross-session
optimum; on Blasco (left, cross-activity) ECG is the signal excluded from the optimum.}
\label{fig:persignal}
\end{figure}

\subsection{Security-bar accounting}
\label{sec:security}
At the equal-error operating point the false-accept rate equals the EER, so we can place the
honest cross-session numbers against published device operating points: Touch ID at FAR
${\approx}\num{2e-5}$ ($\sim$1 in 50{,}000) and Face ID at FAR ${\approx}\num{1e-6}$
($\sim$1 in 1{,}000{,}000). Table~\ref{tab:security} and Figure~\ref{fig:security} give the
gap. The best honest cross-session fusion is on the order of \num{10000} times more permissive
than Touch ID and \num{100000} times more permissive than Face ID: four to five orders of
magnitude short of a real device unlock. Two caveats keep this an order-of-magnitude,
directional statement rather than a same-threshold claim. First, the device FARs are
manufacturer specifications tuned to a low-FAR (and correspondingly high false-reject) point,
whereas our number is the balanced equal-error point; a wearable pushed to a Touch-ID-like FAR
would reject the genuine user a large fraction of the time. Second, our EER is closed-set (the
impostors are the other enrolled subjects) while Touch ID and Face ID are open-set (they may
face a complete stranger); the open-set gap would, if anything, be larger. The exact
multipliers in Table~\ref{tab:security} follow from dividing the EER by the device FAR and are
reported for completeness, not as calibrated system comparisons.

\begin{table}[t]
\centering
\small
\begin{tabular}{llccc}
\toprule
Operating point & Combination & FAR (=EER) & vs.\ Touch ID & vs.\ Face ID \\
\midrule
Blasco cross-activity & PPG+GSR & 0.199 & $\sim$9{,}970$\times$ & $\sim$199{,}400$\times$ \\
Exam-Stress cross-day & GSR+ACC & 0.312 & $\sim$15{,}595$\times$ & $\sim$311{,}900$\times$ \\
\midrule
Touch ID (reference) & n/a & \num{2e-5} & 1$\times$ & n/a \\
Face ID (reference) & n/a & \num{1e-6} & n/a & 1$\times$ \\
\bottomrule
\end{tabular}
\caption{Security-bar accounting for the best \emph{honest} cross-session fusion. ``$\times$''
is how many times more permissive (more false accepts) than the reference.}
\label{tab:security}
\end{table}

\begin{figure}[t]
\centering
\includegraphics[width=0.92\textwidth]{{{artifact:art_e20bdd4d-8584-4c95-8841-3d1d7e54111c}}}
\caption{False-accept rate at the equal-error operating point (log scale). The shaded band is
the device-unlock regime (Face ID / Touch ID). Even the best \emph{within-session} fusion
(0.01) sits well outside it; the honest cross-session points (0.20, 0.31) are near the far
edge.}
\label{fig:security}
\end{figure}

\subsection{Robustness and honesty checks}
\label{sec:robust}
Four controls confirm the reported numbers are genuine, not artefacts of our choices:
\begin{itemize}
  \item \textbf{Split-leakage check.} A random-window within-session split gives EER
  \num{0.014} (Blasco) and \num{0.147} (Exam-Stress); the leak-free time-blocked split gives
  \num{0.030} and \num{0.234}. The difference is temporal-adjacency leakage: autocorrelated
  neighbours straddling a random enrol/probe boundary. All within-session numbers in this paper
  use the leak-free split; the random-split numbers appear only here, to quantify the
  inflation they cause.
  \item \textbf{Label-shuffle (leakage control).} Permuting subject labels within each
  session drives both within- and cross-session EER to chance (Blasco within
  $0.030\!\to\!0.502$, cross $0.210\!\to\!0.497$; Exam-Stress within $0.234\!\to\!0.497$, cross
  $0.339\!\to\!0.497$), confirming the reported performance is genuine identity signal, not
  residual leakage.
  \item \textbf{Window length.} Rebuilding Blasco features at \SI{3}{}, \SI{5}{}, and
  \SI{10}{\second} leaves within-session EER at \num{0.022}-\num{0.038} and cross-session at
  ${\sim}0.21$ throughout, so the cross-session failure is not a windowing artefact.
  \item \textbf{Bootstrap confidence intervals.} The within- and cross-session 95\% CIs do not
  overlap on either dataset (Blasco within $[0.014, 0.047]$ vs.\ cross $[0.119, 0.297]$;
  Exam-Stress within $[0.226, 0.249]$ vs.\ cross $[0.295, 0.377]$). These are resampled over a
  small number of units (3 within-session sessions and 6 cross-session ordered pairs per
  dataset), so they describe variation across sessions rather than subject-level sampling
  uncertainty; we report the unit counts rather than over-claim tight intervals.
\end{itemize}

\section{Discussion}
\label{sec:discussion}

\paragraph{What the number means.} The within-session ${\sim}\SI{2}{\percent}$ EER is real but
answers the wrong question. It measures how well a model separates people \emph{inside a
single sitting}, where session-specific nuisance structure is fully available and fully
predictive. A password lives in the cross-session regime, and there the honest number is
\num{0.20}-\num{0.34} EER: unusable as an unlock.

\paragraph{Why fusion does not save it.} Fusion is often proposed as the fix: combine weak
signals into a strong joint identity. Our data shows a small real gain from fusion, but the
gain is dwarfed by the cross-session failure, and in each dataset a cardiac channel is dropped
from the cross-session optimum: ECG on Blasco, PPG on the wrist device. There is no universal
``cardiac degrades most'' rule (on Blasco the cardiac ECG degrades most across sessions, while
on Exam-Stress the cardiac PPG degrades least because it barely worked within-session in the
first place); the consistent finding is that no channel reaches a usable cross-session
operating point. Fusion cannot rescue a set of signals that individually fail to transfer.

\paragraph{The honest use case.} None of this means wearable signals are useless for identity.
It means the framing must match the physics. A fast fused read cannot be a device unlock, but
it can be a low-stakes \emph{continuity} check (``is this probably still the same wearer as
a moment ago'') where a \SI{20}{\percent} error rate is tolerable because the fallback is
benign (re-prompt, or defer to a real authenticator). Presenting the same technology as a
fingerprint replacement is the error we are trying to forestall.

\paragraph{Limitations.} (1) Blasco's cross-session axis is cross-\emph{activity} on a single
day, a proxy; the true cross-day result comes from Exam-Stress and is if anything worse.
(2) The two datasets' \emph{absolute} EERs are not directly comparable (Exam-Stress wrist
BVP over multi-hour exams is much noisier than Blasco's short controlled recording), so we
compare the within$\to$cross \emph{collapse within each dataset}, not raw numbers across them.
(3) Cross-day, the residual signal is carried by GSR and ACC, which may encode
behavioural/postural habit or per-session stress baseline rather than stable physiological
identity; we therefore do not claim the residual is ``identity'' at all. (4) Both datasets are
small ($n=25$ and $n=10$); the direction of the effect is robust to our controls, but absolute
magnitudes on larger, more diverse populations remain to be measured, which is exactly what
the released benchmark is built to enable.

\section{Conclusion}
Fast multi-signal wearable fusion is strong within a single session (Blasco full-fusion EER
\num{0.030} under a leak-free split) and does not transfer: it degrades about $7\times$ across
activity states on the same day, and on a real wrist device across days it never authenticated
well even within a session (EER \num{0.234}) and stays weak (\num{0.312}). Either way the
honest cross-session operating point is four to five orders of magnitude short of a real device
unlock. Fusion helps a little and does not close the gap; a cardiac channel drops out of the
cross-session optimum in each dataset. The contribution is a reproducible, leakage-audited,
cross-session-honest fusion benchmark that makes this easy to verify and extend, and a plain
statement of the honest use case. We release all code and a one-command reproduction.

\section*{Reproducibility and data availability}
All code, loaders, the leakage audit, the evaluation harness, and \texttt{reproduce.py} are
released at \url{https://github.com/noorabdalla04/wearable-fusion-auth-benchmark}. Raw datasets
are obtained from their original sources under their own licenses (Blasco 2018 and PhysioNet
Exam-Stress); the repository never redistributes raw data and documents provenance and
checksums. Results in this paper were produced with a fixed random seed.

\bibliographystyle{plainnat}
\bibliography{refs}

\end{document}
